\subsection{Every congruence Atkin--Lehner test is trivial}

\paragraph{D23.F18.}
Let $H$ be the marked profinite level-six subgroup of
D23.F6 and let $H^{\rm cong}$ be its full congruence
completion, namely the closure of
$\Gamma_0(6)/\{\pm I\}$ in
$\operatorname{SL}_2(\widehat{\mathbf Z})/\{\pm I\}$.
For every multiplier-one gentle element $h$, its
induced outer action on $H^{\rm cong}$ is trivial.
The three Atkin--Lehner involutions extend to this
completion. Consequently every outer defect
$AL_{6,e}(h)$, $e=2,3,6$, is trivial on the entire
congruence completion and on every finite quotient
of it to which the Atkin--Lehner action descends.

In particular the congruence version of
$AL_{6,3}=1\Longleftrightarrow\Delta=0$ is false:
the supplied E.58 witness has $\Delta=1$ and still
has trivial congruence defect. Any decision of the
full A.86 comparison must use noncongruence
information, not just a finer congruence level.

\begin{proof}
Write $q$ for the adelic congruence representation,
projectivized by the global sign when its source
is $\widehat B$.  The signed covariance theorem
DELTA\_ON\_GT AD.1 gives a uniquely defined
$t_h\in\widehat{\mathbf Z}$ with
\[
 q\Phi_h=\operatorname{Ad}(U^{t_h})q,
 \qquad U=\begin{pmatrix}1&1\\0&1\end{pmatrix}.
\]
The marked cusp generator of $H$ is $a=\tau$,
with $q(a)=U^{-1}$.  Thus $U^t$ belongs to
$H^{\rm cong}$ for every $t\in\widehat{\mathbf Z}$.
This can also be seen directly from its lower-left
entry, which is zero.

For the representative used in D23.F7,
$\psi_h=\operatorname{Inn}(a^{8\rho_h})\Phi_h$,
the congruence covariance is
\[
 q\psi_h
  =\operatorname{Ad}(U^{t_h-8\rho_h})q.
\tag{D23.F18.1}
\]
The conjugating element is itself in $H^{\rm cong}$.
Hence the induced automorphism of that completion
is inner. In particular the congruence kernel of
$H$ is preserved: equation (D23.F18.1) proves this,
so no kernel invariance is assumed in forming
the quotient action.

For precision about bad primes, the closure in
question is exactly the projective image of
\[
 \{g\in\operatorname{SL}_2(\widehat{\mathbf Z}):
                         g_{21}\in6\widehat{\mathbf Z}\}.
\]
Indeed reductions of $\operatorname{SL}_2(\mathbf Z)$
surject onto every finite congruence quotient,
and lifting at a modulus divisible by six retains
the displayed lower-left condition. Thus it
is the inverse image of the usual level-six
subgroup and not an unspecified larger adelic
group.

Every rational Atkin--Lehner matrix normalizing
$\Gamma_0(6)$ acts continuously by conjugation
on its local closures, with a continuous inverse,
and therefore on the full congruence completion.
For the specified $W_3$ this is completely
explicit. If
\[
 W_3=\begin{pmatrix}3&-2\\6&-3\end{pmatrix},
 \qquad g=\begin{pmatrix}a&b\\6c&d\end{pmatrix},
\]
then
\[
 W_3gW_3^{-1}=
 \begin{pmatrix}
 -3a+12c-6b+4d&2a-8c+3b-2d\\
 6(-a+3c-2b+d)&4a-12c+6b-3d
 \end{pmatrix}.
\]
These integral formulas make sense over
$\widehat{\mathbf Z}$, preserve determinant one
and the level-six condition, and their square
is the identity because $W_3^2=-3I$.
The same closure argument applies to $W_2,W_6$;
their classes square to inner automorphisms as
usual for Atkin--Lehner involutions.

In the outer automorphism group of
$H^{\rm cong}$ the class of (D23.F18.1) is
therefore identity. Its crossed commutator with
any Atkin--Lehner class is identity as well.
If a chosen finite normal congruence kernel is
not initially Atkin--Lehner stable, intersect
its finitely many translates under the
Atkin--Lehner group; their squares act inner
and hence fix every normal subgroup. The
result is a finite-index invariant congruence
kernel, and the same inner-action argument
applies there. This explains the precise
finite-quotient formulation of the statement.

Finally E.58.23 and E.65 supply an actual gentle
element with $\Delta=1$. The just-proved
congruence conclusion applies to it, proving
the stated refutation of the congruence version.
No conclusion that its full noncongruence
Atkin--Lehner defect is trivial is drawn.
\end{proof}
